% TOP_PROBLEM: {"id": 3259, "title": "Non-edges vs. feedback edge sets in digraphs", "category_id": 3, "category": {"id": 3, "name": "graph_theory", "display_name": "Graph Theory", "description": "Problems involving graphs, networks, and their properties.", "slug": "graph-theory", "order_index": 3, "created_at": "2026-07-31T15:26:25.671Z"}, "status": "open", "problem_number": "OPG-1793", "set_id": 12}
% FIRST_POSTED: 2026-10-02
% ATTEMPT_STATUS: unresolved
% UnsolvedMath ID 3259; source catalogue code OPG-1793
% !TeX program = lualatex
% GPT-6 Astra Ultra research attempt; independently checked partial work, 2026-10-02.
\documentclass[11pt,a4paper]{article}
\usepackage[margin=25mm]{geometry}
\usepackage{fontspec}
\setmainfont{lmroman10-regular.otf}[BoldFont=lmroman10-bold.otf,
 ItalicFont=lmroman10-italic.otf,BoldItalicFont=lmroman10-bolditalic.otf]
\usepackage{amsmath,amssymb,mathtools}
\usepackage{unicode-math}
\setmathfont{latinmodern-math.otf}
\AtBeginDocument{\let\setminus\smallsetminus}
\usepackage{xurl,hyperref}
\hypersetup{colorlinks=true,urlcolor=blue}
\setcounter{secnumdepth}{0}
\setlength{\parindent}{0pt}
\setlength{\parskip}{5pt}
\setlength{\emergencystretch}{3em}
\begin{document}
\section*{Non-edges vs. feedback edge sets in digraphs}\label{3259}
{\small UnsolvedMath ID 3259 · OPG-1793 · Graph Theory · OpenGarden}
\subsection{Definitions and mathematical statement}
Let $D=(V,A)$ be a finite simple digraph with no directed cycle of length
one, two, or three. Consequently there are no loops or pairs of opposite
arcs. A non-edge is an unordered pair of distinct vertices with no arc
in either direction. Let
\[
 \gamma(D)=\bigl|\{\{u,v\}\subseteq V:uv\notin A, vu\notin A\}\bigr|.
\]
The unordered convention counts each missing adjacency once.

A feedback arc set is a subset $F\subseteq A$ such that $D-F$ has no
directed cycle. Equivalently, $F$ meets the arc set of every directed
cycle of $D$. Define
$\beta(D)=\min\{|F|:F\subseteq A,\ D-F\text{ is acyclic}\}$.
All deletions here are arc deletions, not vertex deletions.

The Chudnovsky--Seymour--Sullivan conjecture asks whether every such
3-free digraph satisfies
\[
                        \beta(D)\le\frac{\gamma(D)}2.
\]
Since $\beta(D)$ is an integer, the right-hand side can equivalently be
replaced by $\lfloor\gamma(D)/2\rfloor$. The condition allows undirected
triangles with a transitive orientation, and permits directed cycles
of length four or more.

One interpretation is that an almost-complete orientation with no short
directed cycle must be close to acyclic. More precisely, a linear ordering
of $V$ has a backward arc when its head precedes its tail. The minimum
number of backward arcs over all such orderings equals $\beta(D)$,
by taking a topological ordering after deleting a minimum feedback
arc set. This gives an equivalent ordering formulation of the target.
\subsection{Short English statement}
In a digraph without directed cycles of length at most three, can all directed cycles be destroyed by deleting at most half as many arcs as there are missing vertex pairs?
\subsection{Research attempt: a cactus case and sharp blow-ups}
\paragraph{Scope and literature check.}
GPT-6 Astra Ultra, 2 October 2026. Missing pairs are unordered and
feedback sets consist of arcs. Chen--Karson--Liu--Shen [S3] prove the
weaker coefficient $0.8616$ for the general conjecture. A September
2026 preprint [S4] establishes a half-coefficient result for a
bipartite, directed-4-cycle-free setting, with its missing pairs
counted between the parts. Those different hypotheses and counting
conventions do not establish the present assertion. The following
argument proves a restricted class and checks that its coefficient
cannot be improved even on dense, strongly connected examples.

\paragraph{1. Reduction to strongly connected pieces.}
Every directed cycle lies in one strongly connected component.
Consequently, if these components are $D_1,\ldots,D_s$, then
$\beta(D)=\sum_i\beta(D_i)$. One inequality follows by restricting
any feedback set, and the other by taking the union of optimal sets
inside the components. Meanwhile $\gamma(D)\ge\sum_i\gamma(D_i)$,
since missing pairs across components contribute nonnegative terms.
Thus a proof for every strongly connected component would suffice.
This reduction does not require deleting or orienting missing pairs.

\paragraph{2. The underlying cactus case.}
Suppose the underlying undirected graph is a cactus, meaning that any
two of its simple cycles share at most one vertex. Each directed cycle
must then traverse one entire undirected cycle block; there are no
other underlying simple cycles on which it could be supported.
Let $t$ be the number of blocks whose orientations are directed cycles.
They are edge-disjoint, so every feedback arc set has at least $t$
arcs. Deleting one arc from each destroys all directed cycles, giving
$\beta(D)=t$.

Every such block has length at least four, by the hypothesis excluding
short directed cycles. An undirected cycle block of length $\ell$
has no chords: a chord would create distinct cycles sharing more than
one vertex. It therefore supplies $\ell(\ell-3)/2\ge2$ missing pairs.
Choose two in each directed block. Different blocks share at most one
vertex, so none of these unordered pairs is counted twice. Hence
$\gamma(D)\ge2t$, proving the proposed inequality for every orientation
of a cactus that satisfies the original cycle restrictions.

\paragraph{3. An equality family with many overlapping cycles.}
Take four blocks $A,B,C,E$, each indexed by $1,\ldots,q$. Orient each
block as a transitive tournament. Between consecutive blocks put all
arcs in the cyclic directions
\[
 A\longrightarrow B\longrightarrow C\longrightarrow E
 \longrightarrow A,
\]
and leave the two opposite block pairs nonadjacent. There is no directed
triangle. A triangle inside a block is transitive; a triangle using
two blocks has uniform cross directions; and three distinct blocks
include a nonadjacent pair. Exactly the opposite block pairs are
missing, so $\gamma=2q^2$.

For every pair $(i,j)$ take the directed cycle
$(A_i,B_j,C_i,E_j,A_i)$. These $q^2$ cycles are arc-disjoint: the two
indices specify each arc uniquely, in all four types of cross arcs.
Thus $\beta\ge q^2$. Conversely delete every arc from $A$ to $B$.
The block order $B,C,E,A$, combined with each internal transitive
order, is a topological ordering of what remains. Exactly $q^2$ arcs
were deleted, so
\[
                 \beta=q^2=\gamma/2.
\]
The case $q=1$ is the directed four-cycle; larger values also show that
sparse cactus structure is not necessary for equality.

\paragraph{Verification and first remaining gap.}
The exact script [S9] constructs these examples for $1\le q\le4$,
checks the triangle restriction, missing-pair count, arc-disjoint
packing, and the asserted acyclic deletion. These are certificates of
the displayed family, not an exhaustive search over digraphs. The
cactus proof assigns two private missing pairs to each necessary
cycle deletion. In general cycles overlap without supplying private
pairs, and the blow-ups illustrate substantial reuse of vertices.
No charging scheme controlling that reuse is obtained here.

\paragraph{Final status.}
UNRESOLVED in the full stated class. The cactus case and the sharp
equality family are proved; neither constitutes a new general bound.
\subsection{Sources}
\begin{itemize}
\item[S1] ULAM.AI and the UnsolvedMath Contributors, supplied problems.json, record 3259 (OPG-1793), OpenGarden collection. Catalogue identity and source transcription; CC BY 4.0.
\url{https://huggingface.co/datasets/ulamai/UnsolvedMath}.
\item[S2] Open Problem Garden, Non-edges vs. feedback edge sets in digraphs. Original problem and discussion; the mathematical formulation below is expanded from the supplied source record.
\url{http://www.openproblemgarden.org/op/non_edges_vs_feedback_edge_sets_in_digraphs}.
\item[S3] Kevin Chen, Sean Karson, Dan Liu and Jian Shen, \emph{On the Chudnovsky--Seymour--Sullivan Conjecture on Cycles in Triangle-free Digraphs} (2009). \url{https://arxiv.org/abs/0909.2468}.
\item[S4] Bin Chen, Jianfeng Hou and Siyue Liu, \emph{Feedback edge set in bipartite digraph} (preprint, September 2026). \url{https://arxiv.org/abs/2609.06462}.
\item[S9] GPT-6 Astra Ultra, exact finite checks accompanying this attempt (2 October 2026). Repository script, Python standard library only. \texttt{scripts/check\_attempt\_batch03\_digraphs\_2026\_10\_02.py}.
\end{itemize}
\end{document}
